\documentclass[11pt,a4paper]{ctexart}
\usepackage[margin=22mm]{geometry}
\usepackage{amsmath,amssymb,booktabs,longtable,hyperref,xcolor}
\hypersetup{colorlinks=true,linkcolor=blue,urlcolor=blue}
\setlength{\parindent}{0pt}
\setlength{\parskip}{4pt}
\newcommand{\code}[1]{\path{#1}}
\newcommand{\fe}[5]{\par\begin{minipage}{\linewidth}\subsubsection*{#1\quad #2}\vspace{-3mm}\[\displaystyle #3\]\textbf{含义/输入输出：}#4\\\textbf{源码与例子：}#5\par\end{minipage}\par\medskip}
\title{CHASI 统一符号、公式与变量去向}
\author{源码核查版；历史编号 C01--C38、L01--L17、Z01--Z03 保留}
\date{2026-10-16}
\begin{document}
\maketitle
\section*{阅读前提与符号}
本文件记录\code{matlab/chasi/chasi_run.m} 当前入口与独立旧链，不把公式当作实际波形发射。所有 $g$ 为资源组，$l$ 为 TRP，$k$ 为当前组内接收 UE，$j$ 为组内求和 UE；全局 UE 编号为 $u=\operatorname{ueIDs}_g(k)$。$N_{\rm RB}=273$，$N_t=64$，$N_r=4$，当前流数 $r_{kg}=2$（可配置且 $r\le\min(N_t,N_r)$）。$H_{klg}\in\mathbb C^{4\times64}$，$w_{ljg}\in\mathbb C^{64\times r_{jg}}$ 已含发射幅度，$A_{kg}\in\mathbb C^{4\times r_{kg}}$，$E_{kg},\Omega_{kg}\in\mathbb C^{r_{kg}\times r_{kg}}$。$p$ 专指由 $w$ 范数导出的功率，不在发射式再乘 $\sqrt p$。$b_{ug}\in\{0,1\}$ 是调度，$a_{lkg}\in\{0,1\}$ 是服务关联。这里的 $g$ 同时作为当前近似的一组数据 RB 与一次 SRS 跳频单元；第17组粒度冲突在末节明列。

代码别名：$w\leftrightarrow$\code{precoder}，$A\leftrightarrow$\code{rxCombiner}，$E\leftrightarrow$\code{mseMatrix}，$\Omega\leftrightarrow$\code{mseWeight}，$\eta\leftrightarrow$\code{userWeight}，$b\leftrightarrow$\code{scheduleMask}，$a\leftrightarrow$\code{servingMask}，$H^{\rm true}\leftrightarrow$\code{H_true}，$\widehat H\leftrightarrow$\code{H}，$C\leftrightarrow$\code{C_stat}。距离专用 $d$，业务掩码用 $\chi$、需求权重用 $\kappa$；固定接收器后的算子记 $\mathcal Q$，避免与旧接收方向 $q^{\rm old}$ 混用。每资源组 $K_g$ 不同；cell 容器不构成新的数学维度。$H^{\rm true}$ 用于观测生成、\code{ideal} 分支和最终评价，非理想分支优化只见 $\widehat H,C$。误差统计采用每发射列相同的 $4\times4$ 接收协方差，等价假设 $I_{64}\otimes C$；实际未生成完整 $256\times256$ 协方差。

\section{拓扑、阵列、传播与测量：C01--C13}
\fe{C01}{七站共址扇区}{s_0=(450,450),\quad s_m=s_0+300(\cos(m\pi/3),\sin(m\pi/3)),\ m=0,\ldots,5}{站址 $s$ 为米；21 TRP 由每站3次复制得到；300 m 是邻站距离。}{\code{chasi_generate_scene>geometry}；例如中心站3个 TRP 坐标相同，主方向相隔120°。项目假设。}
\fe{C02}{UE 扇区位置}{\theta=2\pi((l-1)\bmod3)/3+(U_1-1/2)2\pi/3,\quad d=25+125\sqrt{U_2},\quad q_u=s_l+d(\cos\theta,\sin\theta)}{$U_1,U_2\sim\mathrm{Unif}(0,1)$；$q_u$ 米、$d\in[25,150]$ m；非严格环扇面积均匀。}{\code{chasi_generate_scene>geometry}；$U_2=0.25$ 时 $d=87.5$ m。项目假设。}
\fe{C03}{扇区方向增益}{G_{\rm sec}=0.08+0.92\max(\cos\Delta\theta,0)^4}{方位差 $\Delta\theta$ 弧度，输出无量纲，乘入链路增益。}{\code{chasi_generate_scene>link_gain}；正前方 $G=1$，背面 $G=0.08$。系数为项目假设。}
\fe{C04}{默认归一化增益}{\beta_{ul}=0.025\left(\frac{\max(d_{ul},10)}{100}\right)^{-3.2}G_{\rm sec}}{$d$ 为米，$\beta$ 是默认无量纲 surrogate；不能当标准路径损耗。}{\code{chasi_generate_scene>link_gain}；100m正前方得0.025。项目假设。}
\fe{C05}{可选 UMi/UMa 子集}{\beta_{ul}=\frac{10^{-PL_{ul}/10}G_{\rm sec}}{\sigma_{\rm RB}^2N_{\rm RB}},\quad PL_{ul}=PL_{\rm LOS/NLOS}(d_{3D},f_c)+X_{\rm sh}}{$d_{3D}$ 米、$f_c$ GHz、$PL$ dB、$\sigma_{\rm RB}^2$ W；LOS/NLOS 分支及阴影单独抽样。}{\code{chasi_generate_scene>los_probability,outdoor_pathloss,link_gain}；默认不调用。项目所写标准子集，非完整 CDL/阵列。}
\fe{C06}{小尺度真信道}{R_{\rm rx}[m,n]=0.35^{|m-n|},\quad H^{\rm true}_{ulg}=\sqrt{\beta_{ul}}\,\operatorname{chol}(R_{\rm rx})\Xi_{ulg},\quad \Xi_{mn}\sim\mathcal{CN}(0,1)}{$R_{\rm rx}[4,4]$，$\Xi,H[4,64]$；$H$ 按组独立抽样，未建8×4×2物理阵列。}{\code{chasi_generate_scene} 开头及组循环；例：同一 TRP 的相邻接收行相关系数0.35。}
\fe{C07}{RB与跨组权重}{n_g=16\ (g\le16),\quad n_{17}=17,\quad a_g=n^{\rm tx}_g/273}{数据 RB 总数273；$a_g$ 无量纲，进入目标、功率、速率。}{\code{chasi_config} 与 \code{chasi_matrix_wmmse>rbweight}；完整末组权重17/273。}
\fe{C08}{RB开启规则}{N_{\rm target}=\operatorname{round}(\rho_{\rm PRB}273),\quad \mathrm{active}_{lg}\in\{0,1\}}{$\rho_{\rm PRB}=.4$ 默认；随机整组或共同精确 RB 规则产生组开启。}{\code{chasi_generate_scene>schedule}；.4 给约109个目标RB，与业务30\%不同。}
\fe{C09}{调度与服务}{b_{ug}=\mathbf1\{u\in\mathcal K_g\},\quad a_{lkg}=\mathbf1\{l\in\mathcal S_{kg}\}}{$b[N_{\rm UE},17]$、$a_g[21,K_g]$；每开启小区抽取2/3/4个 UE；服务为 home 加可用最近 TRP。}{\code{chasi_generate_scene>schedule} 与组循环；例如 $b_{u,3}=1$ 且 $a_{1,k,3}=a_{4,k,3}=1$。固定规则。}
\fe{C10}{CSI/SRS噪声与年龄}{C_0=\beta R_{\rm rx},\quad n_c=\frac{\beta}{10^{\gamma_c/10}}e^{2t_c/\tau},\quad n_s=\frac{\beta}{10^{\gamma_s/10}}e^{2t_s/\tau}}{$C_0,C$ 为 $[4,4]$ 归一化协方差；$t_c,t_s,\tau$ ms，$\gamma$ dB。}{\code{chasi_generate_scene} 组循环；10ms更旧的观测增大噪声方差。项目假设。}
\fe{C11}{独立观测后验}{C_{\rm post}=\left(C_0^{-1}+I_4/n_c+S^HS/n_s\right)^{-1},\quad \widehat H=C_{\rm post}\left(Y_c/n_c+S^HY_s/n_s\right)}{$S[2,4]$ 只选两行，$S^HS[4,4]$；$Y_c[4,64],Y_s[2,64]$；未测组省略 SRS 项；输出 $\widehat H[4,64],C_{\rm post}[4,4]$。}{\code{chasi_generate_scene} 的独立模式，\code{calibrate} 用5000列独立样本核查；例：仅第1、2行获得附加 SRS 精度。源码还原。}
\fe{C12}{复用导频条件更新}{S=\sigma_s^2I+\sum_{i\in\mathcal P}\left(\alpha^2C_i+(1-\alpha^2)\beta_iR_{\rm rx}\right)_{\rm rows},\quad \widehat H_i^+=\widehat H_i+\alpha C_i[:,\mathrm{rows}]S^{-1}(Y-\widehat Y)}{$\mathcal P$ 是同导频 UE，$\alpha=e^{-t_s/\tau}$；输出边际 $4\times4$ 后验，未保存跨UE协方差。}{\code{chasi_generate_scene>reuse_update}；两UE复用时观察包含两者叠加，不能按 C11 独立噪声解释。项目高斯假设。}
\fe{C13}{远链截断}{d_{ul}>750\mathrm m\ \land\ a_{lkg}=0\ \Rightarrow\ H_{klg}=0,\quad \mathrm{omittedGain}_{klg}=\beta_{ul}}{真/估计 H 均为零；遗失增益只进诊断，不补回实际干扰。}{\code{chasi_generate_scene} 组循环，\code{chasi_evaluate} 仅报期望遗漏/噪声比。项目假设。}

\section{初始化、发送、接收与目标：C14--C23}
\fe{C14}{初始复预编码}{w^{(0)}_{lkg}=\operatorname{svdRight}_{1:r}(\widehat H_{klg}),\quad w^{(0)}\leftarrow\Pi_{\mathcal C}(w^{(0)})}{初始每块 $[64,r]$，经跨组功率投影后含幅度；服务集外恒零。}{\code{chasi_matrix_wmmse} 第1--39行；例：rank2 取两个右奇异向量。非最终优化。}
\fe{C15}{等效TRP发射}{\mathbb E[s_js_i^H]=\delta_{ij}I,\quad x_{lg}=\sum_{j:a_{ljg}=1}w_{ljg}s_{jg}}{$s_j[r,1]$ 单位协方差，$x_l[64,1]$；不再乘 $\sqrt p$。}{\code{chasi_matrix_wmmse>state,qapply} 只运算二阶量，未生成数据符号或 OFDM 波形。教学信号式。}
\fe{C16}{多TRP到UE}{y_{kg}=\sum_l H_{klg}x_{lg}+n_{kg},\quad F_{kjg}=\sum_lH_{klg}w_{ljg}}{$y[4,1]$，$F[4,r]$；同一 UE 的服务 TRP 先相干求和，不是逐站功率相加。}{\code{chasi_matrix_wmmse>state,trueRates}；两服务TRP时 $F=H_{k1}w_{1j}+H_{k2}w_{2j}$。模型等效。}
\fe{C17}{接收协方差}{p_{lg}=\sum_j\|w_{ljg}\|_F^2,\quad R^{yy}_{kg}=I_4+\sum_jF_{kjg}F_{kjg}^H+\sum_lp_{lg}C_{klg}}{$p_{lg}$ W，$R^{yy}[4,4]$ 为噪声归一化模型；末项只用于统计分支。}{\code{chasi_matrix_wmmse>state}；若某站没有波束则其误差项为零。从源码还原，修复旧 PDF 截断。}
\fe{C18}{MMSE接收矩阵}{A_{kg}=(R^{yy}_{kg})^{-1}F_{kkg}}{$A[4,r]$，接收统计与期望有效信道进入；代码用反斜杠求解。}{\code{chasi_matrix_wmmse>state}；rank2 时4×4系统右侧4×2。}
\fe{C19}{MSE与权重}{E_{kg}=I_r-A^HF_{kkg}-F_{kkg}^HA+A^HR^{yy}A,\quad\Omega_{kg}=E_{kg}^{-1}}{$E,\Omega[r,r]$ Hermitian；与发射波束 $w$ 不同。}{\code{chasi_matrix_wmmse>state}；若 $E=.5I_2$ 则 $\Omega=2I_2$。}
\fe{C20}{组内估计频谱效率}{R^{\rm SE,est}_{kg}=-\log_2\det E_{kg}}{单位 bit/s/Hz；仅在 MMSE 接收及当前高斯模型下与互信息式相联。}{\code{chasi_matrix_wmmse>state}；$E=.5I_2$ 给2 bit/s/Hz。}
\fe{C21}{全局UE组率}{R^{\rm est}_u=\sum_{g:k\mapsto u}a_gR^{\rm SE,est}_{kg}}{$u$ 全局身份，$a_g=n_g^{\rm tx}/273$；输出 $[N_{\rm UE},1]$ bit/s/Hz。}{\code{chasi_matrix_wmmse>objectiveSnapshot}；同一 UE 两组各1 bit/s/Hz时按各组权重相加。}
\fe{C22}{PF及导数权重}{J_{\rm PF}=\sum_{u:b_u=1}\log(1+R^{\rm est}_u/f),\quad\eta_u=\frac{1/(R^{\rm est}_u+f)}{\operatorname{median}_{v:b_v=1}[1/(R^{\rm est}_v+f)]}}{$f=.01$ bit/s/Hz 为 floor；$\eta$ 无量纲公共尺度归一，映射回每组 UE。}{\code{chasi_matrix_wmmse>objectiveSnapshot} PF 分支；小速率UE权重较大，但不保证边缘率提高。项目目标。}
\fe{C23}{非PF目标}{J_{\rm sum}=\sum_g a_g\sum_kR^{\rm SE,est}_{kg},\quad J_{\rm edge}=\sum_ga_g\sum_k\eta^{\rm fixed}_{kg}R^{\rm SE,est}_{kg}}{理想/名义/统计默认分支用求和；edge权重来自 home 大尺度增益且截在[.25,8]。}{\code{chasi_matrix_wmmse>objectiveSnapshot,prepareResources}；两个分支的数值不可当同一目标比较。}

\section{WMMSE 外轮与集中式 ADMM：C24--C31}
以下 $M,B,D,Q$ 是固定 $A,\Omega,\eta$ 后的\textbf{优化系数与算子}，不是新发射、接收信号。函数调用者为 \code{chasi_run}→\code{chasi_matrix_wmmse}；返回 \code{solution.precoder} 给 \code{chasi_evaluate}，其余诊断给保存层。
\fe{C24}{二次系数}{M_{kg}=A_{kg}(\eta_{kg}\Omega_{kg})A_{kg}^H,\quad B_{lkg}=\widehat H_{klg}^HA_{kg}(\eta_{kg}\Omega_{kg}),\quad D_{lg}=\sum_k\operatorname{tr}(M_{kg}C_{klg})}{$M[4,4]$、$B[64,r]$、$D$ 标量；固定接收器和MSE权重，构造预编码子问题。}{\code{chasi_matrix_wmmse>quadratic}：输入 \code{H,C,rxCombiner,mseWeight,userWeights}；如 $C=0$ 则 $D=0$。从源码还原。}
\fe{C25}{隐式二次算子}{[\mathcal Q_g(X)]_{li}=\sum_k\widehat H_{klg}^HM_{kg}\left(\sum_m\widehat H_{kmg}X_{mi}\right)+D_{lg}X_{li}}{$X,w,\mathcal Q(X)$ 同形 `[L,K,64,r]`，非服务块置零；$m$ 是 TRP 求和下标，$i$ 是发射 UE。}{\code{chasi_matrix_wmmse>qapply(H,parts,x,mask)}；两个站的 $X$ 先通过各站 H 求和再左乘 M。旧PDF公式求和下标截断在此修复。}
\fe{C26}{固定接收器的预编码问题}{\min_{w\in\mathcal C}\sum_ga_g\left[\operatorname{Re}\langle w_g,\mathcal Q_g(w_g)\rangle-2\operatorname{Re}\langle B_g,w_g\rangle\right],\quad\mathcal C:\ a_{lkg}=0\Rightarrow w_{lkg}=0,\ \sum_ga_g\sum_j\|w_{ljg}\|_F^2\le P_l}{仅子问题局部二次优化，$P_l=20$ W 默认；$\langle A,B\rangle=\operatorname{tr}(A^HB)$。}{\code{chasi_matrix_wmmse>quadratic,qapply,precoder_admm,project}；例：两组权重16/273与17/273共同耗费同一TRP预算。源码还原；非全局最优保证。}
\fe{C27}{ADMM变量分裂}{\min_{x,z}\ f(x)+I_{\mathcal C}(z)\quad\text{s.t. }x=z,\qquad\mathcal L_\rho=f(x)+I_{\mathcal C}(z)+\frac\rho2\|x-z+y\|_F^2-\frac\rho2\|y\|_F^2}{$x$ 无约束原变量，$z$ 可行副本，$y$ 缩放对偶，$\rho$ 罚参数；均与 $w$ 同形。}{\code{chasi_matrix_wmmse>precoder_admm}；初始化 $x=z=w^{(t)}$、$y=0$。源码还原，旧 PDF 截断修复。}
\fe{C28}{CG求原变量}{\left(a_g\mathcal Q_g+\frac\rho2I\right)x_g=a_gB_g+\frac\rho2(z_g-y_g)}{固定 $A,\Omega,\eta,z,y$；CG 只近似解线性系统，内层最大8次默认；服务外块强制0。}{\code{chasi_matrix_wmmse>precoder_admm,cg}；$\rho=.05$ 默认；\code{cgResidual/cgDirection} 是数值迭代量，不是无线功率。}
\fe{C29}{投影与对偶}{z^{+}=\Pi_{\mathcal C}(x^{+}+y),\quad y^{+}=y+x^{+}-z^{+},\quad z_{lg}=\frac{(x+y)_{lg}}{1+\lambda_l a_g}}{$\lambda_l\ge0$ 用二分使每TRP跨组约束满足，非服务块先置0；返回 $z$，不返回不可行 $x$。}{\code{chasi_matrix_wmmse>project,projectedpower,precoder_admm}；预算未超时 $z=x+y$。旧PDF裁切修复。}
\fe{C30}{内层残差}{r_{\rm pri}=\|x-z\|_F,\quad r_{\rm dual}=\rho\|z-z^{\rm old}\|_F,\quad\epsilon_{\rm pri}=\sqrt n\epsilon_{\rm abs}+\epsilon_{\rm rel}\max(\|x\|,\|z\|)}{另有 $\epsilon_{\rm dual}=\sqrt n\epsilon_{\rm abs}+\epsilon_{\rm rel}\rho\|y\|$；还须CG相对残差达阈值才记收敛。}{\code{chasi_matrix_wmmse>precoder_admm}；12次上限先到则保存可行 $z$ 且 \code{converged=false}。源码还原。}
\fe{C31}{WMMSE回溯与停止}{w^{\rm trial}=w+\alpha(z-w),\quad\alpha\in\{1,1/2,\ldots,2^{-11}\},\quad J(w^{\rm trial})\ge J(w)-10^{-7}\max(1,|J(w)|)}{接受后重算 $A,E,\Omega$ 和目标；外轮还检查目标相对变化、预编码相对变化和最近ADMM收敛；默认2次上限。}{\code{chasi_matrix_wmmse} 第41--103行；若无候选提高则保留旧 $w$ 并报告原因。}

\section{真实信道、业务与输出：C32--C38}
\fe{C32}{真信道组频谱效率}{\begin{aligned}F^{\rm true}_{kj}&=\sum_lH^{\rm true}_{kl}w_{lj},\\R^{\rm SE,true}_k&=\log_2\det\left(I+\sum_jF^{\rm true}_{kj}F^{\rm true,H}_{kj}\right)\\&\quad-\log_2\det\left(I+\sum_{j\ne k}F^{\rm true}_{kj}F^{\rm true,H}_{kj}\right)\end{aligned}}{bit/s/Hz；同一预编码、真实 H、无估计协方差惩罚。}{\code{chasi_evaluate>true_rates} 与 \code{chasi_matrix_wmmse>trueRates}；两用户时另一用户项只入第二个行列式。}
\fe{C33}{RB聚合与容量代理}{R^{\rm physical}_u=\sum_ga_gR^{\rm SE,true}_{kg},\quad c_{ugt}=R^{\rm SE,true}_{kg}n_g^{\rm tx}(12\Delta f)\Delta t}{前者 bit/s/Hz，后者 bit/TTI；$\Delta f=30$ kHz，$\Delta t=1$ ms。}{\code{chasi_evaluate} 主循环；16RB、1 bit/s/Hz、1ms 给5760 bit。}
\fe{C34}{有效忙时体验代理}{N_u^{\rm valid}=\mathbf1\{\exists g:b_{ug}=1\}\sum_t\chi_{ut}^{\rm busy}(1-\chi_{ut}^{\rm terminal}),\quad R_u^{\rm exp}=\frac{\sum_t c_{ut}\kappa_{ut}}{N_u^{\rm valid}\Delta t}}{$\kappa=1$ 普通、$2/7$ 尾段；$R^{\rm exp}$ bit/s 再除 $10^6$ 得Mbps；无有效样本为NaN。}{\code{chasi_evaluate}；末TTI不入分母。容量/业务代理，不是RLC ACK。}
\fe{C35}{分位数与缺样本}{\operatorname{Quantile}_q=x_{(\lfloor v\rfloor+1)}+(v-\lfloor v\rfloor)(x_{(\lceil v\rceil+1)}-x_{(\lfloor v\rfloor+1)}),\quad v=(n-1)q/100}{线性插值经验分位；有效UE集合与无样本置零集合是两个不同分母。}{\code{chasi_evaluate>percentile_or_nan,cdf_candidates}；两个数1、3的5\%分位为1.1。}
\fe{C36}{完整业务样本代理}{R^{\rm burst}=\frac{\sum_{t\in\mathcal T_{\rm eligible}}\mathrm{deliveredBits}_t}{|\mathcal T_{\rm eligible}|\Delta t}}{$\mathcal T_{\rm eligible}$ 排除终止与空缓冲；340 TTI 窗不足完整UE秒；不是已确认协议吞吐。}{\code{chasi_burst_metrics,chasi_second_metrics}；仅含一个终止TTI的 burst 无有效分母。}
\fe{C37}{可选移动队列回放}{q_{t}^{+}=q_t+a_t,\quad D_t^{\rm delivered}=\min(q_t^+,c_t),\quad q_{t+1}=q_t^+-D_t^{\rm delivered}}{队列 bit，$a$ 到达、$c$ 物理容量代理；非正式默认静态求解，回放可独立记录。}{\code{chasi_time_evolution} TTI循环；到达10 bit、容量6 bit，则留下4 bit。}
\fe{C38}{可选时间相关与更新}{\alpha=J_0(2\pi f_D\Delta t),\quad H_{t+\Delta t}=\alpha\sqrt{\beta_{t+\Delta t}/\beta_t}H_t+\sqrt{1-\alpha^2}\,\sqrt{\beta_{t+\Delta t}}L_{\rm rx}Z}{速度3 km/h默认，20 TTI块演化；周期CSI/SRS更新与可选快速SVD波束重置；非在线 WMMSE 重优化。}{\code{chasi_time_evolution}；$\alpha=1$ 时只有大尺度漂移。项目Gauss--Markov假设。}

\section{旧 MATLAB 独立链：L01--L17}
旧链来自 \code{matlab/core/cf_evaluate_candidate.m}，其 $w^{\rm old}$ 的轴序 `[Tx,stream,DU,UE,RBG]`；虽与新 $w$ 同属预编码，但服务/调度、信道、功率约束与率定义不同，\textbf{不存在正式转换接口}。以下仅把旧式改写成统一符号，不表示算法已合并。旧 Q 为预设接收方向，记 $q^{\rm old}$，不能当作新 MMSE $A$。
\fe{L01}{旧合成信道}{H^{\rm old}_{ru g}=\left(\frac{\max(d,d_{\min})}{d_0}\right)^{-\gamma/2}Z\quad(\text{可再全局RMS归一})}{旧 DU/UE/RBG 轴，非 C06 的真实场景 H。}{\code{cf_generate_scenario}；距离加倍时幅度按 $2^{-\gamma/2}$ 缩小。项目假设。}
\fe{L02}{旧SRS误差}{e=\operatorname{clip}\!\left(\frac{m}{1+10^{\mathrm{Presinr}_{\rm dB}/10}}\right),\quad\widehat H=\sqrt{1-e}H+\sqrt e\,\operatorname{RMS}(H)Z}{Presinr dB、缺测乘子 $m$；不是 C11 的两端口矩阵后验。}{\code{cf_apply_srs_measurement_model}；低Presinr使 $e$ 增大。}
\fe{L03}{九维外层候选}{x\in[0,1]^9\mapsto(\beta_{\rm PF},K_{\rm conn},\tau_b,\rho_{\rm link},\rho_{\rm power},r_{\max},\tau_r,p_{\rm repair},K_{\rm repair})}{外层归一化候选，经线性/对数/取整映射；不同于新WMMSE外轮。}{\code{cf_decode_candidate}；第2分量控制候选服务DU个数。}
\fe{L04}{旧关联、rank、方向}{b^{\rm old}_{rug}=\mathbf1\{r\in\mathrm{topK}_{r'}\|\widehat H_{r'ug}\|\},\quad r_{ug}=1+\mathbf1\{\sigma_2^2/(\sigma_1^2+\varepsilon)\ge\tau_r\}}{旧 $b[R,U,G]$ 混合候选关联与调度；$q^{\rm old}$ 取拼接H的左奇异向量。}{\code{cf_evaluate_candidate>buildInitialState}；谱第二值小时rank1。}
\fe{L05}{旧MRT样式初值}{w^{\rm old}_{rusg}=\frac{\widehat H_{rug}^Hq^{\rm old}_{usg}}{\|\widehat H_{rug}^Hq^{\rm old}_{usg}\|+\varepsilon}\quad\text{后续功率缩放}}{先是单位方向，功率缩放及内轮后成为含幅度波束；新 $w$ 无直接赋值。}{\code{cf_evaluate_candidate>buildInitialState,normalizeAllDUPower}；4维接收方向乘 $4\times N_t$ 的H得 $N_t$ 列向量。}
\fe{L06}{旧流级SINR}{z_{us,jt}=h^{\rm eff,H}_{us}w^{\rm old}_{jt},\quad\mathrm{SINR}_{us}=\frac{|z_{us,us}|^2}{\sum_{jt}|z_{us,jt}|^2-|z_{us,us}|^2+\sigma^2+\varepsilon}}{先用 $q^{\rm old}$ 投影为标量；与 C32 的多流logdet不同。}{\code{cf_evaluate_candidate>computeMetrics}；另一UE同RB的波束贡献进分母。}
\fe{L07}{旧调度目标}{\mathrm{WPS}_{us}=\delta_u\ln(1+\mathrm{SINR}_{us}),\quad J_{\rm true}=\sum_{\rm scheduled,u,g,s}\log_2(1+\mathrm{SINR}_{usg})}{WPS 为旧内轮辅助项，不是 MSE 权重；$J_{\rm true}$ 是旧外层主目标。}{\code{cf_evaluate_candidate>computeMetrics,cf_compute_true_objective}。}
\fe{L08}{旧二次变换量}{\zeta_{us}=\frac{\sqrt{\delta_u(1+\mathrm{SINR}_{us})}\,z_{us,us}}{P_{\rm signal}+P_{\rm interf}+\sigma^2+\varepsilon}}{$\zeta$ 是旧标量复辅助量，不是新接收矩阵 $A$。}{\code{cf_evaluate_candidate} 内轮开头；信号幅度为零时 $\zeta=0$。}
\fe{L09}{旧再加权}{\alpha_{rug}=\frac1{p_{rug}+\varepsilon_\alpha},\quad\delta_u=\frac1{(\overline R_u+\varepsilon_{\rm PF})^{\beta_{\rm PF}}},\quad\overline R\leftarrow.85\overline R+.15R}{旧功率稀疏权重、PF历史速率；不同于新 $\eta_u$。}{\code{cf_evaluate_candidate} 内轮；弱率UE的 $\delta$ 增大。}
\fe{L10}{旧闭式波束更新}{T=\sum_{u,s}|\zeta_{us}|^2h^{\rm eff}_{us}h^{\rm eff,H}_{us}+(\mathrm{penalty}+\mathrm{reg})I,\quad w^{\rm old,new}=T^{-1}\left(\zeta^*\sqrt{\delta(1+\mathrm{SINR})}h^{\rm eff}\right)}{这里 $T$ 是旧局部线性系统，非新接收矩阵A；旧代码按DU/RBG更新。}{\code{cf_evaluate_candidate>updateBeamformers}；正则化保证数值可解。旧PDF截断修复。}
\fe{L11}{旧每DU/组功率}{p_{rug}=\sum_s\|w^{\rm old}_{rusg}\|^2,\quad\sum_up_{rug}>P_{\max}\Rightarrow w\leftarrow w\sqrt{P_{\max}/\sum_up_{rug}}}{旧功率约束逐 DU/RBG，非 C26 跨组20W。}{\code{cf_evaluate_candidate>powerFromBeamformer,normalizeAllDUPower}；超限4倍时波束幅度乘1/2。}
\fe{L12}{旧最终关联}{b^{\rm final}_{rug}=\mathbf1\{p_{rug}>\tau_b\}\quad\text{再做可选弱用户修复/归一}}{旧 $b$ 来自功率阈值，与新 $b_{ug}$ 的组调度语义不同。}{\code{cf_evaluate_candidate>repairWeakUsers} 等；功率低于阈值的服务被删除。}
\fe{L13}{旧公平诊断}{\mathrm{Jain}=\frac{(\sum_uR_u)^2}{N_{\rm UE}\sum_uR_u^2+\varepsilon}}{无量纲；5/10分位旧实现为排序索引，非 C35 线性插值。}{\code{cf_evaluate_candidate>computeMetrics}；相同速率给接近1。}
\fe{L14}{旧体验率代理}{R^{\rm exp,old}=\mathrm{ThpVolDl}/\mathrm{ThpTimeDl}}{按旧有效调度、缓冲及尾部样本选择；轴可能为RBG，非正式UE秒。}{\code{cf_compute_experience_rate}；无样本按旧规则置0。}
\fe{L15}{旧目标与评分}{\mathrm{Objective}=\mathrm{proposed.SumRate},\quad\mathrm{Score}=\mathrm{benefit}-\mathrm{cost}-\mathrm{penalty}}{Objective 驱动旧搜索；Score 仅工程诊断/并列，真值重算用同一旧波束。}{\code{cf_evaluate_candidate} 第158--188行与 \code{evaluateOnTrueChannel}。}
\fe{L16}{旧条件原生适配}{H^{\rm native}=H^{\rm old}_{\rm est}/\sqrt{\sigma^2},\quad C^{\rm native}=C^{\rm old}/\sigma^2}{只有旧候选显式启用时将其资源送矩阵求解器；不构成正式 CHASI 场景。}{\code{cf_chasi_native_resources,cf_evaluate_chasi_native_candidate}；输入旧 H 并不等于 C06 的正式 H。}
\fe{L17}{旧外层搜索}{x^{\star}=\operatorname*{arg\,max}_{x\in\mathcal X_{\rm GA/PSO/SRACOS}}\mathrm{Objective}(x)}{粒子速度是 \code{psoVelocity}，与通信预编码无关；外层搜索预算仍按旧评价次数。}{\code{cf_search,cf_sracos_guided_adaptive_gapso_search}；候选目标越高越优先。}

\section{2023参考仓库：Z01--Z03（未进入正式入口）}
原仓库 \code{WMMSE_ADMM_Func.py} 为固定 MR 方向、单流标量模型。此处是源码关系索引；论文原式号尚未核实，不能把下式编号当论文编号。
\fe{Z01}{固定方向功率}{\begin{aligned}\mu_{lk}&=\sqrt{p_{lk}},\quad w^{\rm ref}_{lk}=\mu_{lk}f^{\rm MR}_{lk},\\\mathrm{SINR}_k&=\frac{|\sum_l\mu_{lk}a_{lk}|^2}{\sum_{i,l,m}\mu_{li}B_{k,i,lm}\mu_{mi}+\sigma_k^2}\end{aligned}}{$f^{\rm MR}$ 固定方向，$\mu$ 只调幅度；并非当前自由复矩阵 $w$。}{\code{WMMSE_ADMM_Func.py} 第17--131行；两个AP同相时分子先叠幅度。}
\fe{Z02}{参考标量接收与权重}{v_k=\sqrt{n_k}/D_k^{\rm ref},\quad e_k=1-n_k/D_k^{\rm ref},\quad\omega_k=e_k^{-1}}{$v_k$ 是接收器标量，不是当前旧记号 V 或当前 $w$；$D_k^{\rm ref}$ 是参考式分母，非距离；$\omega$ 是MSE权重。}{同上；$e=.5$ 时 $\omega=2$。}
\fe{Z03}{参考ADMM投影}{q_l=\Pi_{\|\cdot\|_2\le\sqrt{P_l}}(\mu_l+g_l),\quad\mathrm{stop}:\ \|\mu-q\|_F/\|\mu\|_F\le\varepsilon}{每AP欧氏球，$g$ 为对偶；不同于当前跨组加权投影与双残差停止。}{同上；超球半径2倍则投影幅度减半。}

\section*{实际消费者与尚存缺口}
\code{chasi_generate_scene} 生成坐标、H、C、\code{scheduleMask}、\code{servingMask}；\code{chasi_matrix_wmmse} 消费估计H/C并生成 $w,A,E,\Omega,\eta$；\code{chasi_evaluate} 消费 $H^{\rm true},w,$业务掩码生成真实组SE、UE率、体验代理；\code{chasi_run} 写 \code{scene.mat}、\code{*_precoders.mat}、\code{*_per_ue.csv}、\code{summary.json}。本文件公式是源码还原或项目假设，未验证的原图/标准来源请见 \code{CHASI_NATIVE_V2_AUDIT.md}。SRS仅为统计观测模型；末跳16个测量RB与17个数据RB尚未在资源张量上拆开；不存在波形级 SRS/PDSCH/PUSCH、双极化阵列或CPU/TRP真实前传样本。2023论文式号、严谨用户体验协议计数、跨语言同输入随机场景与动态回放仍待验证。
\end{document}
